\subsection{Two actual pointed successes of the rational quintic transfer}
\label{sec:pointed-quintic-success}

\paragraph{Finding A39 (three-primary detectors on an explicit reflection quotient).}
Keep the field $k=\mathbf F_{256}$, the curve $C$, and the
degree-five cover in A36. Write $W_s=\infty$ and
$W_t=(0,0)$ in its reciprocal model. There is a reflection
quotient defined over this same field for which the image $R_s$
of its fixed lift of $W_s$ is Weierstrass, while the image $R_t$
of its fixed lift of $W_t$ has a nonzero three-primary canonical
class. The other two images of the fibre over $W_s$ also have
nonzero three-primary canonical classes. Consequently A29 can
be applied over this stated field to separate each of the two
original Weierstrass points, using the same quintic cover and
an additional cyclic cover of degree three or nine of its
reflection quotient. The resulting etale cover of $C$ has
degree fifteen or forty-five. These are statements about this
fixed example, not a universal pointing theorem.

\begin{proof}
\emph{1. An explicit pointed Kummer realization.}
Use the binary field encoding of A36 and put
\[
 \begin{aligned}
 U(X)&=X^2+80X+224,\qquad V_0(X)=59X+180,\\
 A(X)&=X^5+43X^4+39X^3+106X^2+188,\\
 B(X)&=13X^2+13X+80,\qquad F=A+BY.
 \end{aligned}
 \tag{RPT.1}
\]
The exact polynomial identity in the retained certificate is
\[
 F\,\iota(F)=A^2+ABX^2+B^2(uX^5+X^3+uX)=U^5.
 \tag{RPT.2}
\]
Moreover $F$ vanishes to order five along the degree-two divisor
$D_0$ represented by $(U,V_0)$, and its conjugate is nonzero
there. Indeed $F$ is zero modulo $(U,Y-V_0)$, while
$\gcd(U,BX^2)=1$. Its conjugate $F+BX^2$ is therefore a
unit along $D_0$, so the norm identity assigns all five zeros
there. The polynomial $U$ is separable.
Since $A$ is monic of degree five and $BY$ has smaller pole order
at $W_s$, one obtains
$\operatorname{div}(F)=5D_0-10W_s$.

Thus $D$ can be written $z^5=F$, with deck generator
$\sigma z=12z$, where $12$ has exact order five in $k$.
The class $D_0-2W_s$ has exact order five by A36, so this
Kummer cover is connected and etale. It is the pointed cover
of A36 up to a deck-generator change: $J_C(k)[5]$ has order
five, the field contains $\mu_5$, and the split-fibre
normalization fixes the constant twist.

The reflection is
\[
 \tau(X,Y,z)=(X,Y+X^2,U/z).
 \tag{RPT.3}
\]
Let $P$ be the point above $W_s$ specified by $z/X\to1$.
Let $Q$ be the point above $W_t$ with $z=92=\sqrt{224}$.
The identities $92^5=188$ and (RPT.2) show that $Q$ lies on
$D$. Both $P$ and $Q$ are $k$-rational and fixed by $\tau$.
Put $v:D\to E=D/\langle\tau\rangle$ and
\[
 R_s=v(P),\quad R_t=v(Q),\quad
 R_1=v(\sigma P),\quad R_2=v(\sigma^2P),\quad A_s=R_1+R_2.
 \tag{RPT.4}
\]
The five points above $W_s$ have images $R_s,R_1,R_1,R_2,R_2$.

\emph{2. The canonical divisor and the two tested classes.}
The lower jumps at $W_s,W_t$ are one and three. The completed
local rings used in A33 therefore give
$\operatorname{Diff}(v)=2P+4Q$.
Since $K_C\sim2W_s$ and $\pi$ is etale,
\[
             v^*K_E\sim2\sum_{j=1}^4\sigma^jP-4Q
                         =v^*\bigl(2A_s-2R_t\bigr).
 \tag{RPT.5}
\]
Pullback by $v$ is injective on geometric degree-zero divisor
classes. Indeed if $v^*D_1=\operatorname{div}(g)$, then
$\tau g=cg$ for a constant with $c^2=1$, hence $c=1$ in
characteristic two. The function descends to $E$, and $D_1$
is principal. Consequently $K_E\sim2(A_s-R_t)$.

Write $B_0=A_s-R_t$, a degree-one theta characteristic, and
$Z_R=[R-B_0]$. It is not necessary to identify $B_0$ with the
Raynaud characteristic for this canonical-class computation.
Then
\[
 \begin{aligned}
 \kappa_R&=-2Z_R,\\
 v^*Z_{R_s}&=2P+2Q-\sum_{j=1}^4\sigma^jP,\\
 v^*Z_{R_t}&=4Q-\sum_{j=1}^4\sigma^jP.
 \end{aligned}
 \tag{RPT.6}
\]

\emph{3. The finite Riemann--Roch certificate.}
Here is the complete divisor-space interpretation of the
retained finite matrix. A function in
$\mathcal L(32\pi^*W_s)$ decomposes into its five deck
characters as
\[
           f=\sum_{i=0}^4\frac{z^i}{U^i}
                              \bigl(A_i(X)+B_i(X)Y\bigr),
 \tag{RPT.7}
\]
where
\[
 \deg A_i\le16+i,\quad \deg B_i\le13+i,\quad
 A_i+B_i\bar V_i\equiv0\pmod{U^i}.
 \tag{RPT.8}
\]
The polynomial $\bar V_i$ is the unique curve-equation root
modulo $U^i$ reducing to $V_0+X^2$ modulo $U$.
The congruence removes poles on the conjugate divisor
$\iota D_0$; poles of the coefficient along $D_0$ are canceled
by the zeros of $z^i$. These are all finite places where a
denominator can occur. At infinity $X,Y$ have pole orders two
and five, and their two parities prevent cancellation between
the leading terms of $A_i$ and $B_iY$. This proves precisely
the degree bounds. Conversely multiplication by $U^i$ clears
all finite poles of each deck coefficient, so the affine
coordinate ring $k[X,Y]$ proves that (RPT.7)--(RPT.8) include
every function in this Riemann--Roch space.

For each $i$ an explicit basis consists of
\[
 \begin{gathered}
 (A_i,B_i)=(U^iX^j,0),\quad 0\le j\le16-i,\\
 (A_i,B_i)=
   (\operatorname{rem}_{U^i}(X^j\bar V_i),X^j),
                         \quad0\le j\le13+i.
 \end{gathered}
 \tag{RPT.9}
\]
There are $31$ independent functions for each character,
hence $155$ in total.

At $P$ use the uniformizer $t=Y/X^3$, put $S=1/X=t^2a(t)$,
and write $z=XZ(t)$. Their unique power series are fixed by
\[
 \begin{aligned}
 1+(t+u)a+t^4a^3+ut^8a^5&=0, & a(0)&=u^{-1},\\
 Z^5&=\sum_j A_jS^{5-j}
                   +t\sum_j B_jS^{2-j}, & Z(0)&=1.
 \end{aligned}
 \tag{RPT.10}
\]
In the second line $A_j,B_j$ are the coefficients of (RPT.1).
Both derivatives with respect to the unknown series have
nonzero constant term. At $Q$ use $t=Y$, with
\[
 t^2+X^2t=uX+X^3+uX^5,\quad X(0)=0,\quad
 z^5=A(X)+B(X)t,\quad z(0)=92.
 \tag{RPT.11}
\]
Again the derivatives are units, so these local expansions
are unique. The certificate retains them through degree
$127$ and checks all defining identities to that precision.

A function with divisor
$32(4Q-\sum_{j=1}^4\sigma^jP)$ would belong to (RPT.7), be
regular at $P$, and vanish to order at least $128$ at $Q$.
Regularity at $P$ gives the first $32$ coefficients of
$t^{32}f$ equal to zero; the condition at $Q$ gives another
$128$ coefficients. On the basis (RPT.9) these are the
retained $160$ homogeneous linear conditions in $155$
unknowns. Their matrix has rank $155$ over $\mathbf F_{256}$.
The independent certificate reconstructs the matrix from the
retained residual identities and checks rank $1240$ after
expanding each field entry into its eight-by-eight binary
multiplication matrix. Therefore no such nonzero function
exists, and $32Z_{R_t}\ne0$.

There is also a short positive certificate for $R_s$.
In the expression $\sum_i c_i z^i(A_i+B_iY)/U^i$, take
the following coefficient lists in increasing degree:
\[
\begin{array}{c|c|l|l}
 i&c_i&A_i&B_i\\\hline
 1&225&(224,80,1)&0\\
 2&1&(225,12,162,224)&1\\
 3&13&(224,189,202,174,189)&(13,1)\\
 4&1&(177,0,61,175,79,225)&(236,1,1).
\end{array}
 \tag{RPT.12}
\]
Each summand is regular off $\pi^{-1}(W_s)$ and has pole
order at most two at each point of that fibre. The expansions
(RPT.10)--(RPT.11) give a zero of order four at $P$ and a zero
of order four at $Q$, both with leading coefficient $189$.
The function is nonzero, since its first deck-character
component is $225z$. Its remaining four possible poles have
total degree at most eight, equal to the displayed zero
degree. Thus its divisor is exactly
$4P+4Q-2\sum_{j=1}^4\sigma^jP$.
By (RPT.6) and injectivity of pullback, $2Z_{R_s}=0$ and
$\kappa_{R_s}=0$. Therefore $R_s$ is Weierstrass.

\emph{4. The actual three-primary detectors.}
The rational group has order $2^{10}3^4$ by A36. Its integral
$\mathcal O_5$ action makes its two-primary part a module over
the unramified quadratic ring $\mathcal O_5\otimes\mathbf Z_2$.
At most one nonzero cyclic module summand can occur, since
each contributes two dimensions to the group killed by two
and the geometric two-rank is two. Its order $2^{10}$
therefore gives exponent $32$. Thus $32Z_{R_t}\ne0$ proves
that $Z_{R_t}$, and hence $\kappa_{R_t}$, has a nonzero
three-primary component.

For the other fibre put $Z_i=[R_i-B_0]$ and
$a_0=T(B_0)-2B_0$. The divisor relations are
$T(R_s)=2R_1$ and $T(R_1)=R_s+R_2$.
On the three-primary part $Z_{R_s}=0$, so these relations
give
\[
 a_0=2Z_1,\qquad Z_2=(T+2)Z_1,\qquad
                  Z_{R_t}=Z_1+Z_2=(T+3)Z_1.
 \tag{RPT.13}
\]
The operator $T+3$ is invertible on this part because
$(2-T)(T+3)=5$, and $T+2$ is invertible because
$(1-T)(T+2)=1$. Therefore both $\kappa_{R_1}$ and
$\kappa_{R_2}$ have nonzero three-primary components.

All these points and maps are defined over $k$. Since
$\operatorname{ord}_3(2)=2$, A29 applies to $R_t,R_1,R_2$.
The ring $\mathcal O_5\otimes\mathbf Z_3$ is also unramified
quadratic. A finite module over it of order $3^4$ has exponent
at most nine, so the cyclic detector furnished by A29 has
degree three or nine. Its pullback along the degree-two map
$v$ is connected, since these degrees are coprime, and is
etale over $D$ with the chosen rational lift split.

Finally a source Weierstrass point $W$ needs only one useful
image of a lift. A lifted collision pushed through $v$
becomes equal images by A29. The two possible points in that
quotient fibre are interchanged by $\tau$, whose images on
$C$ are interchanged by the hyperelliptic involution. Since
$W$ is fixed by that involution, both possibilities lie above
$W$. Use $Q$ and $R_t$ for $W_t$, and use $\sigma P$ and
$R_1$ for $W_s$. In both cases the original point collision
is impossible. The etale composite degree is $5\cdot3$ or
$5\cdot9$, as claimed. The composite need not be Galois
over $C$: for $m\in\{3,9\}$ its geometric normal-closure
group embeds in $C_m^5\rtimes C_5$, so has odd order and
the cover belongs to the pro-odd framework.
\end{proof}
